\chapter{Review: Hangul, Sound Changes and Basic Grammar}

KOREA 102R picks up where 101R left off. The first week was a fast review of everything
we were expected to know already: how Hangul syllables are built, the pronunciation rules
that make written and spoken Korean differ, the polite \ko{해요}체 (haeyo-che) present
tense, the core particles, and the two number systems. These notes collect that review in
one place so that the later lessons can refer back to it.

Throughout these notes Korean is followed by its Revised Romanization (RR, the official
South Korean system used on road signs) and an English gloss. RR spells words roughly as
they are \emph{pronounced} after liaison and nasalization, but it does not mark tensed
consonants, so \ko{학교} is written \emph{hakgyo} even though it sounds like
[\ko{학꾜}]. Grammar patterns are written with a tilde, e.g.\ \ko{\til 고 싶다}, where the
tilde stands for the verb stem.

\section{How Hangul syllables are built}

Every Hangul syllable block contains an \textbf{initial consonant} (\ko{초성}, choseong), a
\textbf{vowel} (\ko{중성}, jungseong) and optionally a \textbf{final consonant}
(\ko{받침}, batchim). A syllable cannot start with a vowel in writing, so the silent
consonant \ko{ㅇ} fills the initial slot: \ko{아}, \ko{이}, \ko{우}.

The shape of the block depends on the vowel:
\begin{itemize}
  \item \textbf{Vertical vowels} (\ko{ㅏ ㅑ ㅓ ㅕ ㅣ ㅐ ㅔ}) sit to the \emph{right} of the
        initial consonant: \ko{가}, \ko{너}, \ko{비}.
  \item \textbf{Horizontal vowels} (\ko{ㅗ ㅛ ㅜ ㅠ ㅡ}) sit \emph{below} the initial
        consonant: \ko{고}, \ko{누}, \ko{스}.
  \item \textbf{Compound vowels} (\ko{ㅘ ㅝ ㅙ ㅞ ㅚ ㅟ ㅢ}) wrap around the bottom and
        right: \ko{과}, \ko{뭐}, \ko{의}.
  \item The \textbf{final consonant} always goes along the bottom: \ko{한}, \ko{국},
        \ko{닭} (two final consonants side by side).
\end{itemize}

\begin{center}
\begin{tikzpicture}[x=1cm,y=1cm, every node/.style={font=\small}]
  % C + vertical V
  \begin{scope}[shift={(0,0)}]
    \draw[thick] (0,0) rectangle (2,2);
    \fill[blue!12] (0.05,0.05) rectangle (0.95,1.95);
    \fill[red!12] (1.05,0.05) rectangle (1.95,1.95);
    \node at (0.5,1) {\ko{ㄱ}}; \node at (1.5,1) {\ko{ㅏ}};
    \node at (1,-0.4) {\ko{가} ga};
    \node[font=\scriptsize] at (1,2.3) {C + vertical V};
  \end{scope}
  % C + horizontal V
  \begin{scope}[shift={(2.8,0)}]
    \draw[thick] (0,0) rectangle (2,2);
    \fill[blue!12] (0.05,1.05) rectangle (1.95,1.95);
    \fill[red!12] (0.05,0.05) rectangle (1.95,0.95);
    \node at (1,1.5) {\ko{ㄱ}}; \node at (1,0.5) {\ko{ㅗ}};
    \node at (1,-0.4) {\ko{고} go};
    \node[font=\scriptsize] at (1,2.3) {C + horizontal V};
  \end{scope}
  % C + V + final
  \begin{scope}[shift={(5.6,0)}]
    \draw[thick] (0,0) rectangle (2,2);
    \fill[blue!12] (0.05,1.05) rectangle (0.95,1.95);
    \fill[red!12] (1.05,1.05) rectangle (1.95,1.95);
    \fill[green!15] (0.05,0.05) rectangle (1.95,0.95);
    \node at (0.5,1.5) {\ko{ㅎ}}; \node at (1.5,1.5) {\ko{ㅏ}}; \node at (1,0.5) {\ko{ㄴ}};
    \node at (1,-0.4) {\ko{한} han};
    \node[font=\scriptsize] at (1,2.3) {C + V + final};
  \end{scope}
  % C + horizontal V + final
  \begin{scope}[shift={(8.4,0)}]
    \draw[thick] (0,0) rectangle (2,2);
    \fill[blue!12] (0.05,1.37) rectangle (1.95,1.95);
    \fill[red!12] (0.05,0.71) rectangle (1.95,1.29);
    \fill[green!15] (0.05,0.05) rectangle (1.95,0.63);
    \node at (1,1.66) {\ko{ㄱ}}; \node at (1,1.0) {\ko{ㅜ}}; \node at (1,0.34) {\ko{ㄱ}};
    \node at (1,-0.4) {\ko{국} guk};
    \node[font=\scriptsize] at (1,2.3) {C + horiz.\ V + final};
  \end{scope}
  % compound vowel
  \begin{scope}[shift={(11.2,0)}]
    \draw[thick] (0,0) rectangle (2,2);
    \fill[blue!12] (0.05,1.05) rectangle (1.25,1.95);
    \fill[red!12] (0.05,0.05) rectangle (1.25,0.95);
    \fill[red!25] (1.35,0.05) rectangle (1.95,1.95);
    \node at (0.65,1.5) {\ko{ㄱ}}; \node at (0.65,0.5) {\ko{ㅗ}}; \node at (1.65,1) {\ko{ㅏ}};
    \node at (1,-0.4) {\ko{과} gwa};
    \node[font=\scriptsize] at (1,2.3) {compound V};
  \end{scope}
  \node[font=\scriptsize, blue!60!black] at (2.4,-1.0) {blue: initial};
  \node[font=\scriptsize, red!60!black] at (6.6,-1.0) {red: vowel};
  \node[font=\scriptsize, green!40!black] at (10.6,-1.0) {green: final (\ko{받침})};
\end{tikzpicture}
\end{center}

\begin{remark}
Although there are 27 possible written final consonants (including double ones like
\ko{ㄺ}, \ko{ㅄ}), a syllable-final consonant is only ever \emph{pronounced} as one of
seven sounds: [k], [n], [t], [l], [m], [p], [ng]. For example \ko{부엌} (kitchen) ends in
[k], \ko{옷} (clothes) and \ko{낮} (daytime) end in [t], and \ko{앞} (front) ends in [p].
\end{remark}

\begin{table}[H]
\centering
\begin{tabular}{lll}
\toprule
Pronounced as & Written finals & Examples \\
\midrule
{[k]} & \ko{ㄱ ㅋ ㄲ ㄳ ㄺ} & \ko{책} chaek, \ko{부엌} bueok, \ko{밖} bak \\
{[n]} & \ko{ㄴ ㄵ ㄶ} & \ko{산} san, \ko{앉다} anda \\
{[t]} & \ko{ㄷ ㅌ ㅅ ㅆ ㅈ ㅊ ㅎ} & \ko{옷} ot, \ko{낮} nat, \ko{꽃} kkot \\
{[l]} & \ko{ㄹ ㄼ ㄽ ㄾ ㅀ} & \ko{물} mul, \ko{여덟} yeodeol \\
{[m]} & \ko{ㅁ ㄻ} & \ko{밤} bam, \ko{삶} sam \\
{[p]} & \ko{ㅂ ㅍ ㅄ ㄿ} & \ko{밥} bap, \ko{앞} ap, \ko{없다} eopda \\
{[ng]} & \ko{ㅇ} & \ko{방} bang \\
\bottomrule
\end{tabular}
\end{table}

\section{Sound-change rules}

Written Korean keeps the dictionary spelling of every morpheme; spoken Korean smooths the
joins. The professor insisted we learn these rules \emph{by name}, since the midterm asks
for the pronunciation of a written word in square brackets.

\subsection{Liaison (\ko{연음}, yeoneum)}

When a syllable ending in a consonant is followed by a syllable beginning with the silent
\ko{ㅇ}, the final consonant is pronounced as the initial of the next syllable.

\begin{center}
\begin{tikzpicture}[every node/.style={font=\small}]
  \node[draw, minimum size=1.1cm] (a) at (0,0) {\ko{한}};
  \node[draw, minimum size=1.1cm] (b) at (1.1,0) {\ko{국}};
  \node[draw, minimum size=1.1cm] (c) at (2.2,0) {\ko{어}};
  \draw[-{Stealth}, thick, red!70!black] (1.35,-0.45) to[bend right=50] node[below, font=\scriptsize] {final \ko{ㄱ} moves} (2.0,-0.45);
  \node at (3.6,0) {$\Longrightarrow$};
  \node[draw, minimum size=1.1cm] at (5.0,0) {\ko{한}};
  \node[draw, minimum size=1.1cm] at (6.1,0) {\ko{구}};
  \node[draw, minimum size=1.1cm] at (7.2,0) {\ko{거}};
  \node[font=\scriptsize] at (6.1,-0.9) {[\ko{한구거}] hangugeo, ``Korean language''};
\end{tikzpicture}
\end{center}

\begin{example}[Liaison]
Written form, then [pronunciation]:
\begin{itemize}
  \item \ko{음악} [\ko{으막}] eumak (music)
  \item \ko{읽어요} [\ko{일거요}] ilgeoyo (reads): with a double final, only the
        \emph{second} consonant moves.
  \item \ko{옷이} [\ko{오시}] osi (clothes + subject particle): the \ko{ㅅ} is heard as
        [s] again because it is no longer final.
  \item \ko{좋아요} [\ko{조아요}] joayo (is good): a final \ko{ㅎ} before a vowel
        simply disappears.
  \item Final \ko{ㅇ} never moves: \ko{고양이} [\ko{고양이}] goyangi (cat).
\end{itemize}
\end{example}

\subsection{Nasalization (\ko{비음화}, bieumhwa)}

A stop consonant in final position becomes the nasal made at the same place in the mouth
when the next syllable begins with a nasal (\ko{ㄴ} or \ko{ㅁ}).

\begin{center}
\begin{tikzpicture}[every node/.style={font=\small}, box/.style={draw, rounded corners, minimum width=1.5cm, minimum height=0.7cm}]
  \node[box] (k) at (0,1.6) {\ko{ㄱ} [k]};
  \node[box] (t) at (0,0.8) {\ko{ㄷ} [t]};
  \node[box] (p) at (0,0) {\ko{ㅂ} [p]};
  \node[box, fill=blue!8] (ng) at (5,1.6) {\ko{ㅇ} [ng]};
  \node[box, fill=blue!8] (n) at (5,0.8) {\ko{ㄴ} [n]};
  \node[box, fill=blue!8] (m) at (5,0) {\ko{ㅁ} [m]};
  \draw[-{Stealth}] (k) -- (ng); \draw[-{Stealth}] (t) -- (n); \draw[-{Stealth}] (p) -- (m);
  \node[font=\scriptsize, align=center] at (2.5,2.3) {before \ko{ㄴ} or \ko{ㅁ}};
  \node[font=\scriptsize] at (0,-0.65) {velar / alveolar / bilabial stops};
  \node[font=\scriptsize] at (5,-0.65) {nasals at the same place};
\end{tikzpicture}
\end{center}

\begin{example}[Nasalization]
Each final stop turns into a nasal:
\begin{itemize}
  \item \ko{합니다} [\ko{함니다}] hamnida (does; formal)
  \item \ko{한국말} [\ko{한궁말}] hangungmal (the Korean language)
  \item \ko{작년} [\ko{장년}] jangnyeon (last year)
  \item \ko{끝나요} [\ko{끈나요}] kkeunnayo (it ends): \ko{ㅌ} is first neutralized to
        [t], then nasalized.
  \item \ko{십만} [\ko{심만}] simman (one hundred thousand)
\end{itemize}
\end{example}

Two related rules involve \ko{ㄹ}:
\begin{itemize}
  \item \textbf{\ko{ㄹ} after \ko{ㅁ} or \ko{ㅇ} becomes [n]}: \ko{음료수}
        [\ko{음뇨수}] eumnyosu (soft drink), \ko{정류장} [\ko{정뉴장}] jeongnyujang (bus
        stop).
  \item \textbf{Lateralization} (\ko{유음화}): \ko{ㄴ} next to \ko{ㄹ} becomes [l]:
        \ko{연락} [\ko{열락}] yeollak (contact), \ko{설날} [\ko{설랄}] Seollal (Lunar New
        Year), \ko{신라} [\ko{실라}] Silla (the ancient kingdom).
\end{itemize}

\subsection{Tensification (\ko{경음화}, gyeongeumhwa)}

After a final [k], [t] or [p], a following plain consonant \ko{ㄱ ㄷ ㅂ ㅅ ㅈ} is
pronounced tense (\ko{ㄲ ㄸ ㅃ ㅆ ㅉ}). RR ignores this, so it must be learned by ear.

\begin{example}[Tensification]
\ko{학교} [\ko{학꾜}] hakgyo (school), \ko{식당} [\ko{식땅}] sikdang (restaurant),
\ko{숙제} [\ko{숙쩨}] sukje (homework), \ko{잡지} [\ko{잡찌}] japji (magazine),
\ko{있습니다} [\ko{읻씀니다}] itseumnida (there is; formal).
\end{example}

\subsection{Aspiration (\ko{격음화}, gyeogeumhwa)}

When \ko{ㅎ} meets \ko{ㄱ ㄷ ㅂ ㅈ} (in either order), the two merge into the aspirated
consonant \ko{ㅋ ㅌ ㅍ ㅊ}.

\begin{example}[Aspiration]
\ko{좋다} [\ko{조타}] jota (to be good), \ko{축하해요} [\ko{추카해요}] chukahaeyo
(congratulations), \ko{좋고} [\ko{조코}] joko (good and), \ko{못해요}
[\ko{모태요}] motaeyo (cannot do), \ko{따뜻해요} [\ko{따뜨태요}] ttatteutaeyo (it is
warm).
\end{example}

\begin{important}[Order of operations]
Apply the rules from the \emph{written} form: first neutralize the final to one of the
seven sounds (unless liaison applies), then check the next initial. For \ko{꽃말}
(flower language): \ko{ㅊ} $\to$ [t], then [t] before \ko{ㅁ} $\to$ [n], giving
[\ko{꼰말}] kkonmal. Also watch \textbf{palatalization}: \ko{ㄷ}/\ko{ㅌ} before \ko{이}
become [j]/[ch]: \ko{같이} [\ko{가치}] gachi (together), \ko{굳이} [\ko{구지}] guji.
\end{important}

\section{The polite present: \ko{해요}체}

The \ko{해요} style is the default polite speech level for talking to classmates you do not
know well, shopkeepers, teachers in casual settings, and so on. Every verb and adjective
has a dictionary form ending in \ko{\til 다}; drop \ko{다} to get the \textbf{stem}, then
choose the ending by the stem's last vowel.

\begin{center}
\begin{tikzpicture}[node distance=0.9cm and 1.2cm, every node/.style={font=\small},
    q/.style={draw, diamond, aspect=2.2, inner sep=1pt, align=center},
    r/.style={draw, rounded corners, fill=blue!6, align=center, minimum height=0.8cm}]
  \node[draw, rounded corners] (s) {stem = verb $-$ \ko{다}};
  \node[q, below=of s] (h) {stem is \ko{하}?};
  \node[r, right=of h] (hae) {\ko{해요}\\ \ko{공부해요}};
  \node[q, below=of h] (ao) {last vowel\\ \ko{ㅏ} or \ko{ㅗ}?};
  \node[r, right=of ao] (a) {\ko{\til 아요}\\ \ko{살아요}, \ko{봐요}};
  \node[r, below=of ao] (eo) {\ko{\til 어요}\\ \ko{먹어요}, \ko{마셔요}};
  \draw[-{Stealth}] (s) -- (h);
  \draw[-{Stealth}] (h) -- node[above] {yes} (hae);
  \draw[-{Stealth}] (h) -- node[left] {no} (ao);
  \draw[-{Stealth}] (ao) -- node[above] {yes} (a);
  \draw[-{Stealth}] (ao) -- node[left] {no} (eo);
\end{tikzpicture}
\end{center}

\subsection{Contractions}

When the stem ends in a vowel, the ending contracts with it.

\begin{table}[H]
\centering
\begin{tabular}{lllll}
\toprule
Dictionary & Stem & Rule & \ko{해요} form & Meaning \\
\midrule
\ko{가다} & \ko{가} & \ko{ㅏ}+\ko{아} $\to$ \ko{ㅏ} & \ko{가요} gayo & go \\
\ko{오다} & \ko{오} & \ko{ㅗ}+\ko{아} $\to$ \ko{ㅘ} & \ko{와요} wayo & come \\
\ko{보다} & \ko{보} & \ko{ㅗ}+\ko{아} $\to$ \ko{ㅘ} & \ko{봐요} bwayo & see \\
\ko{서다} & \ko{서} & \ko{ㅓ}+\ko{어} $\to$ \ko{ㅓ} & \ko{서요} seoyo & stand \\
\ko{배우다} & \ko{배우} & \ko{ㅜ}+\ko{어} $\to$ \ko{ㅝ} & \ko{배워요} baewoyo & learn \\
\ko{마시다} & \ko{마시} & \ko{ㅣ}+\ko{어} $\to$ \ko{ㅕ} & \ko{마셔요} masyeoyo & drink \\
\ko{보내다} & \ko{보내} & \ko{ㅐ}+\ko{어} $\to$ \ko{ㅐ} & \ko{보내요} bonaeyo & send \\
\ko{되다} & \ko{되} & \ko{ㅚ}+\ko{어} $\to$ \ko{ㅙ} & \ko{돼요} dwaeyo & become \\
\ko{쓰다} & \ko{쓰} & \ko{ㅡ} drops & \ko{써요} sseoyo & write \\
\bottomrule
\end{tabular}
\end{table}

The copula \ko{이다} (to be) is special: \ko{학생이에요} haksaengieyo (after a consonant)
but \ko{의사예요} uisayeyo (after a vowel). Its negative is \ko{아니에요}: \ko{학생이
아니에요} haksaengi anieyo (I am not a student).

\begin{aside}
\textbf{Speech levels.} Korean verb endings encode how formal and how polite you are
being. In 101R we met the formal \ko{\til (스)ㅂ니다} style (news, presentations,
announcements, the army) and the polite \ko{\til 아요/어요} style. Friends of the same age
and younger siblings get the casual \ko{반말} (banmal), which is the \ko{해요} form with
\ko{요} dropped: \ko{뭐 해?} mwo hae? (what are you doing?). Dropping into \ko{반말} with
someone older without being invited to is rude; a common phrase is \ko{말 놓으세요} mal
noeuseyo, ``please speak casually (to me)''.
\end{aside}

\section{Core particles}

Korean is a subject--object--verb (SOV) language; the verb comes last, and particles
attached to the nouns show their roles, so word order before the verb is fairly free.

\begin{center}
\begin{tikzpicture}[every node/.style={font=\small},
    w/.style={draw, rounded corners, minimum height=0.8cm, fill=gray!8},
    p/.style={draw, rounded corners, minimum height=0.8cm, fill=orange!15}]
  \node[w] (n1) at (0,0) {\ko{저}};
  \node[p, right=0pt of n1] (p1) {\ko{는}};
  \node[w] (n2) at (2.2,0) {\ko{도서관}};
  \node[p, right=0pt of n2] (p2) {\ko{에서}};
  \node[w] (n3) at (5.3,0) {\ko{책}};
  \node[p, right=0pt of n3] (p3) {\ko{을}};
  \node[w, fill=green!12] (v) at (7.6,0) {\ko{읽어요}};
  \node[font=\scriptsize, below=0.25cm of p1] {topic};
  \node[font=\scriptsize, below=0.25cm of p2] {place of action};
  \node[font=\scriptsize, below=0.25cm of p3] {object};
  \node[font=\scriptsize, below=0.25cm of v] {verb (last)};
  \node[font=\scriptsize] at (3.8,1.0) {jeoneun doseogwaneseo chaegeul ilgeoyo. ``I read books at the library.''};
\end{tikzpicture}
\end{center}

\begin{table}[H]
\centering
\begin{tabular}{llll}
\toprule
Particle & After consonant / vowel & Function & Example \\
\midrule
topic & \ko{은} / \ko{는} & what the sentence is about; contrast & \ko{저는 학생이에요.} \\
subject & \ko{이} / \ko{가} & grammatical subject; new information & \ko{비가 와요.} \\
object & \ko{을} / \ko{를} & direct object & \ko{커피를 마셔요.} \\
location (static) & \ko{에} & place of existence; destination; time & \ko{집에 가요.} \\
location (action) & \ko{에서} & place where an action happens & \ko{집에서 쉬어요.} \\
also & \ko{도} & replaces \ko{은/는}, \ko{이/가}, \ko{을/를} & \ko{저도 가요.} \\
\bottomrule
\end{tabular}
\end{table}

\begin{important}[\ko{은/는} versus \ko{이/가}]
\begin{itemize}
  \item Use \ko{은/는} for the topic already known to both speakers, and for
        \emph{contrast}: \ko{커피는 좋아해요. 차는 안 좋아해요.} Keopineun joahaeyo.
        Chaneun an joahaeyo. (Coffee I like; tea I don't.)
  \item Use \ko{이/가} for new information and in answers to ``who/what'' questions:
        \ko{누가 왔어요? --- 민수가 왔어요.} Nuga wasseoyo? --- Minsuga wasseoyo. (Who
        came? --- Minsu came.)
  \item \ko{있다/없다} and adjectives usually take \ko{이/가}: \ko{시간이 없어요.} Sigani
        eopseoyo. (I don't have time.)
  \item Special forms: \ko{나} + \ko{가} $\to$ \ko{내가}, \ko{저} + \ko{가} $\to$
        \ko{제가}, \ko{누구} + \ko{가} $\to$ \ko{누가}.
\end{itemize}
\end{important}

\begin{example}[\ko{에} versus \ko{에서}]
Compare the same room with two different verbs:
\begin{itemize}
  \item \ko{동생이 방에 있어요.} Dongsaengi bange isseoyo. (My younger sibling is in the
        room.) --- existence: \ko{에}.
  \item \ko{동생이 방에서 자요.} Dongsaengi bangeseo jayo. (My younger sibling is sleeping
        in the room.) --- action: \ko{에서}.
  \item \ko{세 시에 학교에 가요.} Se sie hakgyoe gayo. (I go to school at three.) ---
        time and destination: both \ko{에}.
\end{itemize}
\end{example}

\section{Numbers and counters}

Korean has two sets of numbers: \textbf{Sino-Korean} (from Chinese) and \textbf{native
Korean}. Which set to use depends on the counter.

\begin{table}[H]
\centering
\begin{tabular}{rllll}
\toprule
 & Sino-Korean & RR & Native (before counter) & RR \\
\midrule
1 & \ko{일} & il & \ko{하나 (한)} & hana (han) \\
2 & \ko{이} & i & \ko{둘 (두)} & dul (du) \\
3 & \ko{삼} & sam & \ko{셋 (세)} & set (se) \\
4 & \ko{사} & sa & \ko{넷 (네)} & net (ne) \\
5 & \ko{오} & o & \ko{다섯} & daseot \\
6 & \ko{육} & yuk & \ko{여섯} & yeoseot \\
7 & \ko{칠} & chil & \ko{일곱} & ilgop \\
8 & \ko{팔} & pal & \ko{여덟} & yeodeol \\
9 & \ko{구} & gu & \ko{아홉} & ahop \\
10 & \ko{십} & sip & \ko{열} & yeol \\
20 & \ko{이십} & isip & \ko{스물 (스무)} & seumul (seumu) \\
100 & \ko{백} & baek & --- & \\
1000 & \ko{천} & cheon & --- & \\
10\,000 & \ko{만} & man & --- & \\
\bottomrule
\end{tabular}
\end{table}

Native numbers stop being used above about 99 (and most people switch to Sino-Korean well
before that). Note that 10\,000 is its own unit: 50\,000 is \ko{오만} oman (``five
ten-thousands''), and a million is \ko{백만} baengman (``a hundred ten-thousands'').

\begin{table}[H]
\centering
\begin{tabular}{llll}
\toprule
Native counter & Counts & Sino-Korean counter & Counts \\
\midrule
\ko{개} gae & things (general) & \ko{원} won & won (money) \\
\ko{명}/\ko{사람} & people & \ko{년}/\ko{월}/\ko{일} & year / month / day \\
\ko{마리} mari & animals & \ko{분} bun & minutes \\
\ko{권} gwon & books & \ko{층} cheung & floor of a building \\
\ko{병} byeong & bottles & \ko{학년} hangnyeon & school year \\
\ko{잔} jan & cups, glasses & \ko{번} beon & number (bus no.\ 5) \\
\ko{장} jang & flat things (tickets, paper) & \ko{호선} hoseon & subway line \\
\ko{살} sal & years of age & \ko{인분} inbun & portions \\
\ko{시} si & o'clock & \ko{달러} dalleo & dollars \\
\ko{번} beon & times (frequency) & \ko{과} gwa & lesson (textbook) \\
\bottomrule
\end{tabular}
\end{table}

\begin{example}[Counting]
The order is always noun, number, counter:
\begin{itemize}
  \item \ko{사과 세 개} sagwa se gae (three apples).
  \item \ko{학생 스무 명} haksaeng seumu myeong (twenty students).
  \item \ko{지금 두 시 삼십 분이에요.} Jigeum du si samsip bunieyo. (It is 2:30 now.) Hours
        are native, minutes are Sino-Korean.
  \item \ko{이천십팔 년 유월 육 일} icheon sippal nyeon yuwol yuk il (June 6, 2018).
        Note the irregular months \ko{유월} (June, not \ko{육월}) and \ko{시월} (October,
        not \ko{십월}).
\end{itemize}
\end{example}

\begin{aside}
\textbf{Korean age.} When we learned \ko{살} the instructor explained that age has
traditionally been counted from 1 at birth, with everyone adding a year on New Year's Day,
so a Korean may be one or two years ``older'' in Korean age than in international age.
Because age determines speech level and titles (see the lesson on family and honorifics), asking
\ko{몇 살이에요?} myeot sarieyo? or the politer \ko{나이가 어떻게 되세요?} naiga eotteoke
doeseyo? early in a conversation is normal rather than nosy.
\end{aside}

\section*{Review summary}
\begin{itemize}
  \item A syllable block is initial + vowel (+ final); \ko{ㅇ} is a silent placeholder
        initial; finals are pronounced as only seven sounds.
  \item Liaison moves a final onto a following \ko{ㅇ}; stops become nasals before
        \ko{ㄴ/ㅁ}; plain consonants tense after [k t p]; \ko{ㅎ} + stop = aspirated stop.
  \item \ko{해요}체: \ko{하다} $\to$ \ko{해요}; last vowel \ko{ㅏ/ㅗ} $\to$
        \ko{\til 아요}; otherwise \ko{\til 어요}; contract vowels.
  \item Particles: \ko{은/는} topic/contrast, \ko{이/가} subject/new info, \ko{을/를}
        object, \ko{에} static location/destination/time, \ko{에서} place of action.
  \item Native numbers with \ko{개, 명, 마리, 권, 시, 살}; Sino-Korean with \ko{원, 분,
        년, 월, 일, 층}.
\end{itemize}
